% Source-only follow-on fragment. Analytical statements are not Lean acceptance.
% Uses the existing imliliu2024 citation key (published CPC version).

\subsection{The next cyclic boundary}

Cycle rank three is not a single theta case.  Its structural possibilities
can nevertheless be listed exactly.  In the following statement,
suppression replaces a maximal path whose internal vertices have degree
two by one edge; parallel edges are retained.

\begin{proposition}[The four rank-three cores]
\label{density-rank-three-core-list}
Let $J$ be a finite simple $2$-connected bipartite graph with
$\beta(J)=|E(J)|-|V(J)|+1=3$.  Suppressing its degree-two vertices gives
one of the following loopless multigraphs:
\[
 B_4,\qquad T_{2,2,1},\qquad C_4^{\mathrm{alt}},\qquad K_4.
\]
Here $B_4$ has two vertices and four parallel edges;
$T_{2,2,1}$ is a triangle with two doubled sides;
and $C_4^{\mathrm{alt}}$ is a four-cycle with two opposite sides doubled.
Every edge of the suppressed core is replaced in $J$ by a positive-length
path, and these paths have disjoint interiors.  Bipartiteness is equivalent
to the existence of branch-vertex colours $\sigma:V(K)\to\{0,1\}$ such that
\[
 L_{uv}\equiv\sigma(u)+\sigma(v)\pmod2
\]
for every replacement path.  Original simplicity requires at most one
length-one path in each parallel family.
\end{proposition}
\begin{proof}
Write $n=|V(J)|$ and $m=|E(J)|$.  Two-connectivity gives minimum degree
two.  As $J-v$ remains connected, $m-\deg(v)\ge n-2$, so
$\deg(v)\le m-n+2=4$.  The degree excess is
\[
 \sum_v(\deg(v)-2)=2m-2n=4.
\]
Thus the branch degrees are $(4,4)$, $(4,3,3)$ or $(3,3,3,3)$.
There are at least two branch vertices.  A suppressed loop would leave
its interior disconnected from the other branch vertices after deleting
its endpoint, and is therefore impossible.

For two branch vertices the four edges are parallel.  For three vertices,
the three multiplicities satisfy $a+b=4$, $a+c=3$ and $b+c=3$, giving
$(a,b,c)=(2,2,1)$.  With four cubic vertices, a simple kernel is $K_4$.
If there is a doubled pair, a third parallel edge would disconnect that
pair from the other vertices.  The remaining stubs must meet distinct
vertices: if both met the same vertex, the last vertex could not attain
degree three.  The remaining degrees then force the opposite doubled
pair, giving $C_4^{\mathrm{alt}}$.

Alternating colours along each replacement path proves the parity
criterion in both directions.  Only parallel length-one paths can create
duplicate original edges.  Conversely the four kernels have no bridge
or cut vertex, so every simple positive-length subdivision satisfying
the parity condition is two-connected.  Deleting an internal path vertex
uses connectivity after removing its kernel edge; deleting a branch
vertex uses connectivity after removing that branch vertex.  Suppression
and subdivision preserve $m-n+1$.
\end{proof}

All subdivisions of $B_4$ allowed by this proposition are generalized
theta graphs.  Their four lengths have the same parity, and both cases
are Sidorenko~\citep[Theorem~1.3 and Proposition~4.3]{imliliu2024}.
Hence a hypothetical non-Sidorenko single atom of rank three must have
one of the other three suppressed cores.  This is an exclusion, not the
construction of a counterexample.  The cited paper also proves several
uniform or symmetric subdivision cases, but does not supply the
unrestricted nonuniform conclusion needed to raise the present cutoff
from two to three.

\paragraph{A rank-three conjecture.}
For every finite simple $2$-connected bipartite graph $J$ of cycle rank
three and every symmetric $[0,1]$-valued kernel $W$ on a finite probability
space, one may ask whether
\[
 t(J,W)\ge (\mathbb E W)^{|E(J)|}.
\]
The classification makes this a question about three residual
multigraph cores and their admissible path lengths, rather than an
unspecified search over obligatory systems.  It is a working conjecture
here, not a claim that a new open problem has been identified.  Our
bounded literature check has neither proved it nor certified that every
residual family is unresolved in the full literature.

For the higher-uniformity paper, an even weaker question is sufficient:
for every $r\ge3$ and symmetric $r$-kernel $U$, does the same inequality
hold when $W=M_U$ is its two-coordinate marginal?  Indeed,
\[
 t(J^{[r]},U)=t(J,M_U),\qquad \mathbb E M_U=\mathbb E U.
\]
The unrestricted graph-kernel inequality implies this atom inequality.
The converse cannot be inferred by treating every graph kernel as an
available marginal.

\begin{proposition}[A restriction on the marginal class]
\label{density-binary-marginal-restriction}
Let $r\ge3$ and let both states of $\Omega=\{0,1\}$ have positive
probability.  If a symmetric nonnegative $r$-kernel $U$ satisfies
$M_U(0,0)=M_U(1,1)=0$, then $U=0$ pointwise.
\end{proposition}
\begin{proof}
Each diagonal marginal is a finite sum of nonnegative values
$U(a,a,z)$ with strictly positive coefficients.  Its vanishing forces
every such value to vanish.  Every binary $r$-tuple has two equal
coordinates; symmetry moves these to the first two positions.
\end{proof}

Thus a nonzero binary graph kernel with zero diagonal is not an
$r$-kernel marginal on this same carrier.  The full-support,
nonnegativity, symmetry and $r\ge3$ conditions are essential.  This
observation limits one reduction route; it does not disprove any
rank-three density conjecture or rule out other converse arguments.

The useful story is therefore precise: low cyclic complexity supplies
positive cases, while the next rank is reduced to a short list of
joint-kernel inequalities.  A proof must control those joint kernels,
not only the scalar densities of individual paths or pieces.
